\clearpage\subsection{Implementation}
\begingroup
\titlespacing*{\subsubsection}{0pt}{9pt}{4pt}
\setlength{\intextsep}{7pt}
\captionsetup{skip=4pt}
\renewcommand{\arraystretch}{1.08}
The implementation below is a proposed starting configuration. Named libraries, schemas and model settings will be explored during the feasibility stage.

\subsubsection{Application structure}
Next.js/React and TypeScript will use \href{https://tanstack.com/query/latest/docs/framework/react/overview}{TanStack Query} for fetching and job polling. FastAPI/Pydantic will validate JSON contracts \cite{nextjs,fastapi}. SQLAlchemy/Alembic will manage database access/migrations. A Python worker will collect data, fit models and run scenarios.



\subsubsection{API and frontend integration}
\begin{table}[H]
\smalltable
\caption{Proposed endpoints and compact Pydantic contracts.}
\begin{tabularx}{\linewidth}{@{}L{47mm}Y Y@{}}
\toprule\tablehead
\textbf{Endpoint} & \textbf{Input schema} & \textbf{Response schema} \\
\midrule
\texttt{GET /events} & asset, as\_of, cursor & events[], next\_cursor \\
\texttt{GET /wallets/\{address\}} & event\_id, as\_of & score, coverage, activity[] \\
\texttt{POST /analyses} & AnalysisRequest (below) & 202: job\_id, status \\
\texttt{GET /jobs/\{id\}} & job UUID & status, result?, error? \\
\texttt{GET /runs/\{id\}} & run UUID & metrics, units, as\_of, evidence[] \\
\texttt{POST /assistant/query} & session\_id, question, run\_id? & 202: job\_id, status \\
\bottomrule
\end{tabularx}
\end{table}
\textbf{AnalysisRequest:} event\_id: UUID, as\_of: UTC timestamp, horizon: $1|5|20$, holdings: [\{asset: string, weight: float\}]. Weights must be non-negative and sum to one. Invalid inputs return 422 with field errors. AssistantAnswer contains text, run IDs and citations. Cache keys include input and data/model versions. Stale or incomplete evidence is flagged.

\subsubsection{Data ingestion and storage}
HTTPX collectors will fetch incremental pages. Polars/PyArrow will preserve normalised Parquet evidence. Checkpoints, retries and unique source keys will support safe replay. Fills, transfers, splits, merges and redemptions will be reconciled, with sampled Polygon receipt checks.

\begin{table}[H]
\smalltable
\caption{Core physical schema. PK = primary key, FK = foreign key. Internal IDs use UUIDs.}
\begin{tabularx}{\linewidth}{@{}L{31mm}Y@{}}
\toprule\tablehead
\textbf{Table} & \textbf{Key fields and constraints} \\
\midrule
market\_snapshot & id PK, market\_id FK, observed\_at/collected\_at timestamps, price numeric(18,8), evidence\_uri text, source\_key text unique. \\
wallet\_score & id PK, address text, event\_id FK, cut\_off timestamp, score float, history\_count integer, version text. \\
analysis\_run & id PK, event\_id FK, request JSONB, snapshot\_ids JSONB, model\_version FK, seed integer, results JSONB. \\
model\_version & id text PK, estimator text, params JSONB, train\_cut\_off timestamp, artifact\_uri text. \\
job & id PK, run\_id FK nullable, kind/status text, attempts integer, heartbeat timestamp, error JSONB. \\
evidence\_chunk & id PK, event\_id FK, source\_uri text, body text, available\_at timestamp, embedding vector(384), embedding\_version text. \\
assistant\_session & id PK, owner\_id FK, summary text, context JSONB, updated\_at timestamp. \\
\bottomrule
\end{tabularx}
\end{table}
Market/time and wallet/event/cut-off indexes support dated lookups. Numeric results stay in SQL/Parquet. Embeddings support document retrieval only.

\clearpage
\subsubsection{Analytical pipeline}
\begin{table}[H]
\smalltable
\caption{Example feature spaces. Each row is constructed at an information cut-off $t$.}
\begin{tabularx}{\linewidth}{@{}L{28mm}Y L{38mm}@{}}
\toprule\tablehead
\textbf{Task / one row} & \textbf{Candidate features $X_t$} & \textbf{Target / output} \\
\midrule
Volatility / asset-date & Past 20-day mean $r^2$, raw event probability, adjusted minus raw probability & Future $RV_{t,H}$ for $H=1,5,20$ \\
Probability / event-date & Raw probability, spread, anonymous directional flow, eligible-wallet weighted flow & Resolved YES/NO label, then $p_{\mathrm{adj}}$ \\
Clustering / wallet-date & Turnover/capital, net YES--NO exposure, event concentration, typical holding duration & Behaviour group, with no outcome label \\
\bottomrule
\end{tabularx}
\end{table}

\begin{figure}[H]
\centering
% Generated synthetic values. Seed 42. No empirical research claims.
\def\regpoints{(0.2000,0.4631) (0.3130,0.2195) (0.4261,0.7284) (0.5391,0.8531) (0.6522,0.2383) (0.7652,0.4731) (0.8783,0.8955) (0.9913,0.8680) (1.1043,1.0190) (1.2174,0.8974) (1.3304,1.3924) (1.4435,1.4471) (1.5565,1.3554) (1.6696,1.6892) (1.7826,1.6100) (1.8957,1.3707) (2.0087,1.7446) (2.1217,1.5051) (2.2348,2.0252) (2.3478,1.8815) (2.4609,1.9282) (2.5739,1.8883) (2.6870,2.4243) (2.8000,2.1729)}
\def\regintercept{0.18668}
\def\regslope{0.73917}
\def\toyvar{3.93619}
\def\toyes{4.57566}
\expandafter\def\csname cluster0\endcsname{(0.1900,-0.6423) (0.2573,-0.5561) (0.2489,-0.5483) (0.3699,-0.6488) (0.1841,-0.6977) (0.2631,-0.4645) (0.2120,-0.7008) (0.1623,-0.5219) (0.2720,-0.5348) (0.1734,-0.5721) (0.2282,-0.5738) (0.2810,-0.5732) (0.2675,-0.5919) (0.2402,-0.5242) (0.1180,-0.6384) (0.1871,-0.6767)}
\expandafter\def\csname cluster1\endcsname{(0.7780,0.7294) (0.7307,0.6662) (0.6654,0.5098) (0.8130,0.6203) (0.8569,0.6452) (0.7721,0.4945) (0.8686,0.5270) (0.6979,0.4140) (0.7264,0.6097) (0.8114,0.6329) (0.7658,0.5690) (0.8500,0.5129) (0.8365,0.4706) (0.7710,0.5042) (0.7043,0.6084) (0.7624,0.5515)}
\expandafter\def\csname cluster2\endcsname{(0.5137,0.1336) (0.5266,0.0682) (0.4504,0.0704) (0.3619,-0.0937) (0.3874,-0.0397) (0.5080,-0.0287) (0.4535,0.2359) (0.4551,0.1685) (0.4146,0.0553) (0.4135,0.0393) (0.5388,-0.1273) (0.5104,0.1085) (0.4384,-0.0935) (0.4850,0.0165) (0.4963,0.0826) (0.5921,0.0513)}
\def\clustercentres{(0.2284,-0.5916) (0.7757,0.5666) (0.4716,0.0405)}

\begin{tikzpicture}
\begin{axis}[name=reg,width=5.2cm,height=4.4cm,scale only axis=false,
 title={\textbf{Regression}},title style={font=\small\sffamily},
 xlabel={Past 20-day mean $r^2$},ylabel={Next 5-day $RV$},
 label style={font=\scriptsize},tick label style={font=\scriptsize},
 xmin=0,xmax=3,ymin=0,ymax=3,xtick={0,1,2,3},ytick={0,1,2,3},
 axis lines=left,grid=major,grid style={gray!15},clip=false]
\addplot[only marks,mark=*,mark size=1.3pt,gray] coordinates {\regpoints};
\addplot[black,thick,domain=0:3]{\regintercept+\regslope*x};
\node[font=\tiny,anchor=north west] at (rel axis cs:0,1) {Units: $10^{-4}$};
\end{axis}
\begin{axis}[name=clust,at={(reg.east)},anchor=west,xshift=1.2cm,width=5.2cm,height=4.4cm,
 title={\textbf{Wallet clustering}},title style={font=\small\sffamily},
 xlabel={Turnover / capital per day},ylabel={Directional exposure},
 label style={font=\scriptsize},tick label style={font=\scriptsize},
 xmin=0,xmax=1,ymin=-1,ymax=1,xtick={0,.5,1},ytick={-1,0,1},axis lines=left,grid=major,grid style={gray!15}]
\addplot[only marks,mark=o,mark size=1.5pt,black] coordinates {\csname cluster0\endcsname};
\addplot[only marks,mark=triangle*,mark size=1.5pt,gray] coordinates {\csname cluster1\endcsname};
\addplot[only marks,mark=square*,mark size=1.4pt,black] coordinates {\csname cluster2\endcsname};
\addplot[only marks,mark=x,mark size=4pt,very thick,black] coordinates {\clustercentres};
\end{axis}
\begin{axis}[at={(clust.east)},anchor=west,xshift=1.2cm,width=5.2cm,height=4.4cm,
 title={\textbf{Scenario loss}},title style={font=\small\sffamily},
 xlabel={Portfolio loss (\%)},ylabel={Density},label style={font=\scriptsize},tick label style={font=\scriptsize},
 xmin=-3,xmax=6,ymin=0,ymax=.34,xtick={-2,0,2,4,6},ytick={0,.15,.3},axis lines=left,
 declare function={mix(\x)=.6*exp(-((\x-2)/1.4)^2/2)/(1.4*sqrt(2*pi))+.4*exp(-((\x+.5)/.8)^2/2)/(.8*sqrt(2*pi));}]
\addplot[draw=none,fill=gray!40,domain=\toyvar:6,samples=40]{mix(x)}\closedcycle;
\addplot[black,thick,domain=-3:6,samples=100]{mix(x)};
\draw[dashed] (axis cs:\toyvar,0)--(axis cs:\toyvar,.26);
\node[font=\tiny,anchor=east] at (axis cs:\toyvar,.28) {95\% VaR};
\draw[dotted] (axis cs:\toyes,0)--(axis cs:\toyes,.17);
\node[font=\tiny,anchor=west] at (axis cs:\toyes,.19) {ES};
\end{axis}
\end{tikzpicture}

\caption{Synthetic examples, not research results. A fitted line, wallet groups with centroid crosses, and a loss mixture with its worst 5\% shaded. ES is the mean loss in that tail.}
\end{figure}

\begin{table}[H]
\smalltable
\caption{Candidate implementations in \href{https://scikit-learn.org/stable/supervised_learning.html}{scikit-learn}. The team will explore all five techniques and retain models supported by validation.}
\begin{tabularx}{\linewidth}{@{}L{32mm}L{49mm}Y@{}}
\toprule\tablehead
\textbf{Technique} & \textbf{Estimator / role} & \textbf{Initial exploration} \\
\midrule
Linear regression & LinearRegression / variance & Recent-return baseline and event-augmented fit. \\
Logistic regression & LogisticRegression / probability & Regularisation strength and wallet-feature contribution. \\
SVM & SVC / outcome, SVR / variance & Linear or RBF kernel, $C$, gamma. Calibrate classifier scores chronologically. \\
MLP & MLPClassifier / MLPRegressor & One small hidden layer, e.g.\ 8 or 16 units, L2 penalty. \\
Clustering & StandardScaler + KMeans & Explore $k=2$--$5$, separation and stability across dated samples. \\
\bottomrule
\end{tabularx}
\end{table}

\begin{figure}[H]
\centering
\begin{tikzpicture}[x=1cm,y=1cm,b/.style={flowbox,text width=4.25cm,font=\small\sffamily,inner sep=5pt,minimum height=1.35cm}]
\node[b] (a) at (0,0) {\textbf{Historical event windows}\\Joint asset returns\\grouped by YES / NO};
\node[b] (b) at (5.2,0) {\textbf{NumPy sampler}\\YES: $p$ / NO: $1-p$\\Apply portfolio weights};
\node[b] (c) at (10.4,0) {\textbf{Scenario results}\\Loss probability\\Volatility / VaR / ES};
\draw[flowarrow] (a)--(b);\draw[flowarrow] (b)--(c);
\end{tikzpicture}
\caption{Scenario execution preserves joint asset returns. The seed and evidence manifest are saved.}
\end{figure}
\textbf{Experiment controls.} Pipeline will fit scaling/imputation on training folds only. Chronological, event-grouped validation will purge overlapping targets, including for MLP tuning. Scores, clusters and probability fits use past data only. The regression baseline stays fixed. Clusters do not establish skill or ownership. Incomplete positions exclude affected features. Insufficient scenario evidence triggers an unavailable result or a labelled, reviewed stress assumption.

\clearpage
\subsubsection{Research assistant integration}
\href{https://docs.crewai.com/en/concepts/flows}{CrewAI Flows} is the proposed orchestrator. A typed Pydantic state will hold session\_id, event\_id, assets, horizon, as\_of, evidence IDs and run IDs. Three agent roles will work inside an explicit retrieve--calculate--explain flow. The workflow graph controls execution. Event--asset relationships will initially use relational joins, with knowledge-graph retrieval reserved for later exploration.

\begin{figure}[H]
\centering
\begin{tikzpicture}[x=1cm,y=1cm,b/.style={flowbox,text width=4.25cm,font=\small\sffamily,minimum height=1.45cm,inner sep=5pt}]
\node[b] (r) at (0,0) {\textbf{Event research agent}\\Find evidence\\Propose asset mappings};
\node[b] (a) at (5.2,0) {\textbf{Analysis coordinator}\\Validate inputs\\Request an analysis run};
\node[b] (e) at (10.4,0) {\textbf{Explanation agent}\\Read saved metrics\\Cite evidence sources};
\node[b] (rag) at (0,-1.9) {\textbf{RAG retrieval tool}\\Event and time filter\\Semantic retrieval};
\node[b] (calc) at (5.2,-1.9) {\textbf{Calculation tools}\\Python services\\Optional MCP adapter};
\node[b] (res) at (10.4,-1.9) {\textbf{Evidence and run store}\\Source chunks\\Structured results};
\node[b] (emb) at (0,-3.8) {\textbf{Document indexing}\\Chunk and embed\\Store in pgvector};
\node[b,text width=9.45cm] (mem) at (7.8,-3.8) {\textbf{Application-managed context and memory}\\Session summary and selected inputs\\Saved run references scoped to each user};
\draw[flowarrow] (r)--(a);\draw[flowarrow] (a)--(e);
\draw[flowarrow] (emb)--(rag);\draw[flowarrow] (rag)--(r);
\draw[flowarrow] (a)--(calc);\draw[flowarrow] (calc)--(res);\draw[flowarrow] (res)--(e);
\draw[flowarrow] (mem.north -| a.south)--(calc.south);
\end{tikzpicture}
\caption{Proposed assistant workflow. Retrieved prose provides context. Typed calculation outputs provide numerical claims. New event--asset mappings require review before use.}
\end{figure}

\begin{table}[H]
\smalltable
\caption{Concrete retrieval, tool and memory design choices.}
\begin{tabularx}{\linewidth}{@{}L{27mm}Y@{}}
\toprule\tablehead
\textbf{Mechanism} & \textbf{Proposed implementation} \\
\midrule
RAG + embeddings & Archive event rules, official releases and reviewed mappings. Split into about 180 model-token chunks with 30-token overlap. Sentence Transformers \href{https://huggingface.co/sentence-transformers/all-MiniLM-L6-v2}{all-MiniLM-L6-v2} produces 384-dimensional embeddings. Store source, date and embedding version with every chunk. \\
Retrieval & \href{https://github.com/pgvector/pgvector}{pgvector} cosine search plus PostgreSQL keyword search. Merge ranked matches and return an initial top five. Filter by event, user access and information cut-off before ranking. Exact prices and positions come from SQL tools. \\
Function / tool use & CrewAI \href{https://docs.crewai.com/en/concepts/tools}{BaseTool} with Pydantic arguments: search\_evidence(query, event\_id, as\_of), get\_wallet\_evidence(address, event\_id, as\_of), run\_analysis(AnalysisRequest), get\_run(run\_id). Tools return evidence IDs, job IDs or typed results. \\
MCP option & A \href{https://docs.crewai.com/en/mcp/overview}{CrewAI MCP adapter} may expose the same approved functions to other clients. MCP is a tool interface, not a memory store. Direct Python tools are sufficient for the initial single deployment. \\
Context + memory & Load selected inputs, the latest six messages, a session summary, retrieved chunks and referenced results into a bounded prompt. Store owner-scoped summaries/preferences in PostgreSQL. Reload numerical facts from run IDs. Summaries never become research evidence. \\
\bottomrule
\end{tabularx}
\end{table}

\subsubsection{Deployment and reproducibility}
Docker Compose will run frontend, API, worker and PostgreSQL/pgvector services. A PostgreSQL job queue with retries and heartbeat expiry will recover interrupted work. Record prompts, tool traces, embedding/model versions and source manifests. Keep authenticated access, secrets, persistent volumes and backups. The assistant will stop after bounded tool calls, report missing evidence and preserve dashboard access when the hosted model fails.

\endgroup
